% Generated draft: requires mathematical and source review.
\documentclass[11pt]{article}
\usepackage[margin=25mm]{geometry}
\usepackage{fontspec}
\setmainfont{Latin Modern Roman}
\usepackage{amsmath,amssymb,mathtools,unicode-math}
\setmathfont{Latin Modern Math}
\usepackage{enumitem,needspace,xurl,hyperref}
\setlength{\parindent}{0pt}
\setlength{\parskip}{5pt}
\setlength{\emergencystretch}{3em}
\newcommand{\N}{\mathbb{N}}
\newcommand{\Z}{\mathbb{Z}}
\newcommand{\Q}{\mathbb{Q}}
\newcommand{\R}{\mathbb{R}}
\newcommand{\C}{\mathbb{C}}
\newcommand{\F}{\mathbb{F}}
\newcommand{\A}{\mathbb{A}}
\newcommand{\PP}{\mathbb{P}}
\newcommand{\E}{\mathbb{E}}
\newcommand{\eps}{\varepsilon}
\newcommand{\dd}{\,\mathrm{d}}
\begin{document}
\section*{UnsolvedMath ID: 20000738}
\textbf{Source problem code:} catalogue entry
\subsection{Definitions and mathematical statement}
This permanent record is an aggregated extraction of the whole document \emph{Problem sessions: ``Subconvexity bounds for $L$-functions''}, G. Ricotta, version December 8, 2006, including its contents and Monday--Friday sessions. It cannot be identified with one mathematical problem despite metadata number ``1''. All assertions about what was known, suggested, or in progress below are historical statements of those notes, not a present-day literature assessment. The extraction damages inequalities, subscripts, roots and some hypotheses. The following retains the separate agenda and then states the accompanying aid's distinct conductor lemma.

The Monday unknown-cases discussion asks, following Reznikov, for convexity bounds in every aspect for Rankin--Selberg $L(\pi\times\pi',s)$ with cuspidal representations of $\operatorname{GL}_m$ and $\operatorname{GL}_n$, $m,n\geq1$, and for a geometric method beyond the cited $\operatorname{GL}_2\times\operatorname{GL}_2$ $L^2$ argument. Michel's proposed subconvex targets ask for some $\delta>0$ and every $\varepsilon>0$:
\[
 \begin{aligned}
 |L(f,1/2+it_f)|&\ll_\varepsilon t_f^{1/4-\delta+\varepsilon},&(1.1)\\
 |L(f\times g,1/2+it_f)|&\ll_{g,\varepsilon}t_f^{1/2-\delta+\varepsilon},&(1.2)\\
 |L(f\times g,1/2+it)|&\ll_{f,g,\varepsilon}t^{1-\delta+\varepsilon},&(1.3)\\
 |L(\operatorname{Sym}^2 f,1/2+it)|&\ll_{f,\varepsilon}t^{3/4-\delta+\varepsilon},&(1.4)\\
 |L(\operatorname{Sym}^2 f,1/2+it)|&\ll_{q_f,t,\varepsilon}t_f^{1/2-\delta+\varepsilon}.&(1.5)
 \end{aligned}
\]
Here $f$ is level-one Hecke--Maass in (1.1); $g$ is fixed in (1.2); both forms are fixed in (1.3); and the symmetric-square discussion allows level $q_f$, spectral parameter $t_f$ and nebentypus $\chi_f$. The source's printed Laplace-eigenvalue line is corrupt and is not adopted as an eigenvalue definition. It asks about possible work of Luo and applications of (1.1), noting its conductor is of order $t_f$. For (1.2) it discusses Phillips--Sarnak deformation/nonvanishing when $g$ is holomorphic of weight four and level $q$, including its historical positive-proportion assertion. Jutila suggests extra spectral averaging for (1.3); the extracted Jutila--Motohashi bounds contain the expressions $t^{1+\varepsilon}/\sqrt{t_f}$ and $(t+t_f)^{2/3+\varepsilon}$, with range boundaries involving $t_f^{3/2}$ and $t_f^{2-\varepsilon}$. Their broken typography requires verification before a precise uniform theorem can be quoted; the question is to extend the range.

The symmetric-square level target (1.6) asks for a saving from $q_f^{1/2}$ for trivial or nonquadratic $\chi_f$, and from $q_f^{1/4}$ otherwise, with a factor $q_f^\varepsilon$. Its printed implied-constant subscript also contains $q_f$, so the intended level-uniformity is unresolved. The notes give $L(f\times f,s)=L(\chi_f,s)L(\operatorname{Sym}^2 f,s)$, credit Michel with a Cauchy--Schwarz connection to unequal-form Rankin--Selberg subconvexity, and suggest metaplectic $\widetilde{\operatorname{GL}}_2$ Eisenstein integral representations.

For triple products the same session asks for savings from the following scales:
\[
 \begin{array}{ll}
 k_h^{2-\delta+\varepsilon}&\text{for }L(f\times g\times h,1/2+it),\ h\text{ holomorphic, fixed }f,g,t,q_h,\quad(1.7),\\
 t^{2-\delta+\varepsilon}&\text{for the same value with fixed }f,g,h,\quad(1.8),\\
 t_f^{1-\delta+\varepsilon}&\text{for }L(f\times f\times h,1/2),\text{ fixed }q_f,h,\quad(1.9),\\
 t_f^{2-\delta+\varepsilon}&\text{for }L(f\times g\times h,1/2)\text{ with comparable, unequal,}\quad(1.10)\!\!\\
 &\text{non-conductor-dropping spectral parameters, as suggested by Reznikov.}
 \end{array}
\]
The $k_h$ interpretation of the damaged weight symbol in (1.7) follows its stated holomorphic-$h$ context and needs checking. Merely unequal parameters are not a quantitative nonresonance condition; uniform separation is not specified.

The Monday universality discussion asks why Weyl's conductor exponent $\tfrac14(1-\tfrac13)$ and Burgess' $\tfrac14(1-\tfrac14)$ recur. It lists Weyl-type examples in $\operatorname{GL}_2$ height/spectral aspects, same-level twists and Rankin--Selberg spectral aspects; Burgess-type examples in Dirichlet level and character-conductor twists. It seeks examples and applications improving those exponents. Soundararajan suggests conductor $3^n$ and Vinogradov methods in the depth aspect; Graham--Ringrose methods for suitably factorable squarefree twist moduli are another suggestion. Michel cites Edixhoven's conditional Andr\'e--Oort application for curves in $X_0(1)^2$, but the extracted projection hypothesis and small-split-prime logarithmic cutoff are damaged. The GRH-for-imaginary-quadratic-fields qualification is retained. Fouvry's cited Linnik--Vinogradov consequence concerns growing counts of split primes below $|d|^{1/4+\varepsilon}$; improved character-sum bounds might improve that range. Duke--Friedlander--Iwaniec's class-group amplification problem is also cited. None of this is asserted as a current theorem in the extracted form.

Tuesday asks for period definitions and examples. Lindenstrauss proposes, with $G$ over $\Q$ and subgroup $H$, pairings of the form
\[
 \int_{H(\Q)\backslash H(\A_{\Q})} f(h)g(h)\,dh.
\]
Examples are unipotent Fourier coefficients on $\operatorname{GL}_2$, torus periods linked to central values up to gamma factors, and diagonal triple-product periods linked to $L(f_1\times f_2\times f_3,1/2)$. Quotients by centers, convergence, measures and test-vector normalizations are not fully specified; the extracted real unipotent integral should not be treated as a verified convergent formula. The session asks for a general harmonic-analysis notion of convexity and Lindel\"of for periods, for Miller's comparison of geometric and automorphic periods including extra $L$-value factors, and for a canonical test-vector-independent normalization. Small-group topics additionally include conductor dropping, improved Weyl exponents, Harcos' suggested spectral decomposition of shifted convolutions, explicit Good--Motohashi identities, associativity identities in conductor aspect, period/representation-theory formulations and quantitative equidistribution.

Wednesday includes Silberman's general hybrid target
\[
 |L(\pi,1/2+it)|\ll_\varepsilon Q(\pi,1/2+it)^{1/4-\delta+\varepsilon}
\]
for some $\delta>0$, and discussion of Diaconu--Garrett--Goldfeld's extension of Good's method to $\operatorname{GL}_n\times\operatorname{GL}_{n-1}$. Venkatesh requests applications, including higher-rank equidistribution; Duke mentions higher-rank Artin $L$-functions. Further questions concern period analogues of convexity/Lindel\"of, and Silberman's spectral-aspect sup-norm prediction. Jutila asks for omega results and applications for second-moment errors; the extraction prints $\operatorname{Error}(T)=\Omega(\texttt{pT})$, with an unrecoverable radical/symbol, so no exact exponent is assigned. Venkatesh asks when first moments over families with nearly constant conductor $Q(\mathcal F)$ admit asymptotics, cites the heuristic size comparison $4\log|\mathcal F|>\log Q(\mathcal F)$ and the historical Conrey--Iwaniec cubic-moment benchmark $6\log|\mathcal F|=\log Q(\mathcal F)$, and asks for improvement. These are not universal family theorems. Soundararajan asks about subconvexity off the critical line, especially near its edge, mentioning Vinogradov for zeta near real part one. Michel asks for a $p$-adic analogue, with integrality suggested in connection with Mazur and Prasad.

Thursday discusses Venkatesh's use of quantitative equidistribution for power-saving moment asymptotics, including $\operatorname{GL}_2$ height and twist-height second moments; the extracted latter integral lacks its test weight. Higher-rank obstacles are explicit period representations and archimedean calculations, then ergodic or spectral analysis, with $O_n\times O_{n-1}$ and $\operatorname{GL}_n\times\operatorname{GL}_{n-1}$ singled out. Garrett's Diaconu--Garrett--Goldfeld example is cited. Suggested applications relate $\operatorname{GL}_n$ subconvexity to zeta's $2n$-th moment, including Motohashi-type $\operatorname{GL}_3$/sixth-moment links and Conrey--Iwaniec.

Friday proposes a reference/attempts website and emphasizes higher-rank methods. The historical discussion names Venkatesh--Lindenstrauss, Bernstein--Reznikoff, Ratner and Clozel ideas, and Diaconu--Garrett--Goldfeld/Good work in progress, as reported by Garrett (whose spelling varies in the extraction). Venkatesh suggests fixed $\operatorname{GL}_3$ times varying $\operatorname{GL}_2$, first testing an Eisenstein specialization through
\[
 \sum_{\substack{f\ \mathrm{level\ one\ Hecke\!\! -\!\! Maass}\\t_f\asymp T}}
 \int_{t\asymp T}|L(f,1/2+it)|^6\,dt.
\]
The document describes family size about $T^3$ and conductor size about $T^{12}$, without enough normalization details to reconcile or upgrade this diagnostic to a theorem. It asks where the $\operatorname{GL}_3$ theory enters the analytic analysis.

Separately, the accompanying aid states a finite archimedean conductor lemma for an $r$-fold tempered spherical $\operatorname{GL}_2$ tensor product. Suppose
\[
 t_j=a_jT+O(1),\quad a_j\geq0,\qquad \tau=bT+O(1),
\]
with fixed slopes, and use the conductor convention, up to bounded parity shifts,
\[
 C_\infty(T)\asymp\prod_{\epsilon\in\{\pm1\}^r}
       \left(1+\left|\tau+\sum_{j=1}^r\epsilon_jt_j\right|\right).
\]
If $N=\#\{\epsilon:b+\sum_j\epsilon_ja_j=0\}$ counts resonances with multiplicity, the aid reports
\[
 C_\infty(T)\asymp T^{2^r-N}.
\]
Two-sided comparability, rather than a unit leading constant, is appropriate for the stated $O(1)$ perturbations. If exactly $k$ slopes are positive, the sharp count is
\[
 N\leq2^{r-k}\binom{k}{\lfloor k/2\rfloor}.
\]
A multiscale version sums individual gamma-factor exponents: if each factor is $T^{\rho_\epsilon+o(1)}$, the conductor is $T^{\sum_\epsilon\rho_\epsilon+o(1)}$. These are archimedean conductor statements, not subconvex estimates, and do not resolve the many independent historical questions above.
\subsection{Short English statement}
This is a damaged extraction of an entire 2006 workshop agenda, not one well-defined problem. Its questions concern subconvexity, periods, moments and equidistribution. The supplied aid isolates a conductor-collision counting lemma; it does not solve that agenda. Historical claims and damaged formulas remain explicitly qualified.
\subsection{Sources}
\begin{itemize}
\item Local UnsolvedMath registry, permanent ID $20000738$, source code \texttt{catalogue entry}; complete supplied document extraction and research summary. Canonical local record: \texttt{aim-analytic-number-theory-notes.json}, \texttt{notes[101]}.
\item G. Ricotta, \emph{Problem sessions: ``Subconvexity bounds for $L$-functions''}, version December 8, 2006, AIM workshop \emph{Subconvexity bounds for L-functions}, supplied URL \url{https://aimath.org/WWN/subconvexity/subconvexity.pdf}. No PDF was retrieved; the supplied extraction alone was used.
\item Attributions and named work retained from that historical document include Reznikov, Michel, Luo, Phillips--Sarnak, Jutila--Motohashi, Soundararajan, Graham--Ringrose, Edixhoven, Fouvry, Linnik--Vinogradov, Duke--Friedlander--Iwaniec, Lindenstrauss, Miller, Harcos, Silberman, Diaconu--Garrett--Goldfeld, Good, Venkatesh, Duke, Conrey--Iwaniec, Mazur, Prasad, Bernstein--Reznikoff, Ratner and Clozel. The extraction supplies no full bibliographic entries for these references.
\item Research attempt 1 is labelled invalid statement, problem not a problem, in the background; registry classification is \texttt{OPEN-TRIAGE}, status open, and statement reconstructed/unverified. Its separate conductor lemma is supplied with rigor label ``proved'' and novelty confidence ``low''. These metadata are not independent mathematical review. The machine-generated aid is not peer reviewed and requires independent verification before citation.
\end{itemize}
\end{document}
